\chapter*{Interludium: Was ein Horizont wirklich ist}
\addcontentsline{toc}{chapter}{Interludium: Was ein Horizont wirklich ist}
\markboth{INTERLUDIUM}{INTERLUDIUM}

\begin{oberflaeche}
Halten wir kurz an, ganz ohne Formeln.

Ein Horizont ist keine Länge und kein Format. Er ist der Raum, in dem
eine Zahl gerade lebt -- mit einer festen, endlichen Anzahl von
Plätzen, die sich aus den Fakultäten ergibt: zwei, sechs,
vierundzwanzig, hundertzwanzig. Eine gewöhnliche Zahl kennt nur
ihren Wert. Eine Hori-Zahl kennt ihren Wert \emph{und} ihren Raum,
und beides gehört zu ihr wie bei einem Buch der Text und das Regal,
in dem es steht.

Das klingt nach einer Kleinigkeit, ist aber ein Perspektivwechsel.
Unsere übliche Mathematik nimmt sich den unendlichen Zahlenstrahl
als gegeben und stellt alle Zahlen gleichzeitig darauf. Die
Horimetrik baut stattdessen Treppen: Jede Stufe ist endlich und
vollständig überschaubar, und die Unendlichkeit steckt nicht in
einer einzelnen Stufe, sondern im Versprechen, dass es immer eine
nächste gibt. Nichts an der klassischen Unendlichkeit ist damit
widerlegt -- sie wird nur anders erzählt: als Folge von Räumen statt
als ein einziger Raum.

Der Rest dieses Buchs untersucht, was dieses Erzählen kostet und was
es einbringt. Es kostet Bequemlichkeit: Beim Umzug von Stufe zu
Stufe ändern alle Ziffern ihre Bedeutung, und man muss sich
entscheiden, was beim Umzug erhalten bleiben soll -- der Wert oder
die Gestalt. Es bringt Struktur ein: Genau aus dieser Entscheidung
entstehen die zwei Identitäten, die zwei Rechenwelten und der Satz,
dass man nicht in beiden gleichzeitig leben kann.
\end{oberflaeche}
